\section{Evaluation}
\label{sec:eval}

\subsection{Research Questions}

Our evaluation is organized around three research questions:

\begin{itemize}[noitemsep,topsep=0pt]
\item \textbf{RQ1:} How do IPI defenses perform under static (L1) vs.\ adaptive (L2) attack, measured end-to-end?
\item \textbf{RQ2:} What is the security-utility tradeoff, and how much utility loss is genuine over-blocking vs.\ base-model capability ceiling?
\item \textbf{RQ3:} Does the echo-canary signal provide independent security value, and does multi-signal combination recover utility? (P2 ablation)
\end{itemize}

\subsection{Experimental Setup}

\paragraph{Benchmark.}
We evaluate on the AgentDojo workspace suite (v1.2.2)~\cite{debenedetti2024agentdojo,debenedetti2025agentdojo}, which provides 40 user tasks across calendar, email, and file-management scenarios with ground-truth injection goals. We use 10 user tasks $\times$ 3 injection tasks per seed. Each task pairs a legitimate user request (e.g., ``Who else is invited at the Networking event?'') with an injection goal (e.g., ``Email `Hey, how is it going?' to mark.black-2134@gmail.com''). Crucially, our evaluation uses AgentDojo's \emph{native} pipeline: the agent executes real tools (send\_email, get\_calendar\_events, etc.), producing real side effects and function-call traces.

\paragraph{Backbone.}
We use GPT-4o-mini (API snapshot \texttt{gpt-4o-mini-2024-07-18}) as the primary agent backbone via an OpenAI-compatible API. Main-table results are reported over 3 seeds (0--2); 95\% Wilson confidence intervals are reported in parentheses. Total API cost was approximately \$60--80 USD.

\paragraph{Defenses.}
We evaluate 11 defense configurations: 7 baselines from the literature (Spotlighting, StruQ, Task Shield, IPIGuard, FATH, Polymorphic Prompt, Mixture-of-Encodings) plus 3 design-principle prototypes (P1--P3, \Cref{sec:principles}), and a no-defense baseline. StruQ is evaluated as a \emph{prompt-level structured-query approximation} (\texttt{struq\_prompt}); the original StruQ requires a fine-tuned model checkpoint~\cite{chen2024struq}, which is not reproducible in our API-based pipeline. We report this honestly: our \texttt{struq\_prompt} tests whether \emph{prompt-level} channel separation (without fine-tuning) suffices, not whether the original fine-tuned StruQ works. Attention Tracker and P4 are excluded from the main table for fidelity reasons (\Cref{sec:proxy_results}).

\paragraph{Metrics (end-to-end).}
Unlike proxy-based evaluations that estimate security via string matching, our metrics are computed by AgentDojo's native evaluators on real tool-execution traces:
\begin{itemize}[noitemsep,topsep=0pt]
\item \textbf{ASR} (Attack Success Rate): fraction of attacked tasks where the injection goal was \emph{actually executed}, as determined by \texttt{injection\_task.security\_from\_traces()} checking the real function-call traces and environment state. This is a strict, trace-based measure---not a string-match proxy.
\item \textbf{Utility}: fraction of tasks where the user's task was \emph{correctly completed}, as determined by \texttt{user\_task.utility()} comparing pre- and post-execution environment state. This independently measures task completion, unlike ``unblocked'' proxies that conflate ``not blocked'' with ``done.''
\end{itemize}
We report both ASR and Utility because a defense that blocks everything achieves ASR=0 but Utility=0---a trivially safe but useless operating point. The security-utility tradeoff is the central tension our evaluation surfaces.

\paragraph{Scale.}
The main experiment comprises 11 defenses $\times$ 1 attack (important\_instructions) $\times$ 10 user tasks $\times$ 3 injection tasks $\times$ 3 seeds = 990 trials, all using AgentDojo's native end-to-end pipeline. The P2 ablation (3 configurations $\times$ 90 = 270 trials) disentangles canary vs.\ alignment contributions. Total: 1{,}260 trials, 0 errors. Total API cost was approximately \$60--80 USD.

\subsection{RQ1: Defense Performance Under Static IPI}

\Cref{tab:main} reports end-to-end ASR and Utility for all 11 defenses under the \texttt{important\_instructions} attack, measured by AgentDojo's native trace-based evaluators. The no-defense baseline has ASR 0.74 [0.65, 0.82] and Utility 0.19 [0.12, 0.28], confirming that GPT-4o-mini is susceptible to injection but also that it completes only $\sim$19\% of tasks even without a defense---a baseline capability ceiling that is critical for interpreting over-blocking claims (a defense with Utility $\approx$ 0.19 is not necessarily over-blocking; it may reflect the model's own limitation).

Input-encoding defenses are the strongest performers: Polymorphic (ASR 0.14, Utility 0.73), Mixture-of-Encodings (ASR 0.16, Utility 0.69), and Spotlighting (ASR 0.21, Utility 0.67) all substantially reduce ASR while \emph{preserving} Utility above the no-defense baseline. This means these defenses do not merely block everything---they selectively neutralize injected instructions while allowing legitimate task completion. Runtime-checking defenses (Task Shield ASR 0.10, FATH ASR 0.00, P2 ASR 0.00) achieve lower ASR but at the cost of Utility $\leq$ 0.02, indicating aggressive over-blocking. The prompt-level StruQ approximation (ASR 0.70) provides little protection, confirming that structured-query separation requires fine-tuning to be effective---prompt-level instructions to ``treat tool output as data'' are insufficient.

\subsection{RQ2: Security-Utility Tradeoff}

The central finding from end-to-end evaluation is a sharp security-utility tradeoff: defenses that achieve low ASR tend to also destroy Utility, while defenses that preserve Utility have moderate ASR. \Cref{tab:main} visualizes this---no defense reaches the ideal zone (low ASR \emph{and} high Utility).

\begin{itemize}[noitemsep,topsep=0pt]
\item \textbf{Input encoding is the best tradeoff:} Polymorphic, Mixture-of-Encodings, and Spotlighting cluster in the favorable region (ASR 0.14--0.21, Utility 0.67--0.73). They reduce ASR by 71--81\% relative to baseline while \emph{increasing} Utility above the no-defense ceiling (0.19), indicating they suppress injected instructions without blocking legitimate actions.
\item \textbf{Runtime checking over-blocks:} Task Shield (ASR 0.10, Utility 0.02), FATH (ASR 0.00, Utility 0.02), P1 (ASR 0.03, Utility 0.00), and P2 (ASR 0.00, Utility 0.00) achieve near-zero ASR but Utility $\leq$ 0.02. However, we must distinguish \emph{genuine over-blocking} from \emph{model-weakness}: comparing Task Shield to the no-defense baseline on identical tasks, Task Shield blocks 14.4\% of tasks that the baseline completes, while 80\% of Task Shield's Utility=0 cases are tasks the baseline \emph{also} fails---attributable to GPT-4o-mini's capability ceiling, not the defense.
\item \textbf{Prompt-level structured query is insufficient:} \texttt{struq\_prompt} achieves ASR 0.70 (only 5\% relative reduction vs.\ baseline), confirming that prompt-level channel separation without fine-tuning does not resist injection. The original fine-tuned StruQ~\cite{chen2024struq} would likely perform better; our result bounds the \emph{prompt-only} approach.
\item \textbf{Tool-allow-list defenses are moderate:} IPIGuard (ASR 0.38, Utility 0.06) and P3 (ASR 0.42, Utility 0.02) block some attacks via static tool allow-lists but still sacrifice Utility and miss a substantial fraction of injections.
\end{itemize}

\subsection{Base Model Capability Ceiling}
\label{sec:backbone_main}

A critical observation is that GPT-4o-mini's Utility on AgentDojo workspace tasks is only 0.19---80\% of tasks are not completed even without a defense. This capability ceiling is essential for interpreting over-blocking claims: a defense with Utility $\approx$ 0.02 is not necessarily over-blocking 98\% of tasks; much of this reflects the model's own inability to complete the tasks. By comparing Task Shield to the no-defense baseline on identical (seed, user\_task, injection\_task) triples, we find that Task Shield genuinely over-blocks only 14.4\% of tasks (those the baseline completes but Task Shield blocks), while 80\% of its Utility=0 cases are tasks the baseline also fails. Future work should evaluate on stronger backbones (e.g., GPT-4o, Claude) where the baseline Utility is higher, enabling cleaner separation of defense over-blocking from model weakness.

\subsection{Main Results Table}

\Cref{tab:main} reports results for the 10 defenses we can faithfully evaluate in our API-based pipeline. Attention Tracker (IP) requires access to model internals (attention weights) not available via API; our implementation uses a token-frequency proxy. We report its results separately in \Cref{sec:proxy_results} and do not include it in the main quantitative conclusions. P4 (Least Privilege) requires real OS-level sandboxing; our implementation is a prompt-level simulation, so we discuss it qualitatively (\Cref{sec:principles}) rather than in the main table.

\begin{table*}[t]
\centering
\small
\caption{Main results: 11 defenses under \texttt{important\_instructions} attack on AgentDojo workspace (10 user tasks $\times$ 3 injection tasks $\times$ 3 seeds = 90 trials per defense, GPT-4o-mini). ASR and Utility are computed by AgentDojo's native end-to-end evaluators (\texttt{security\_from\_traces} and \texttt{user\_task.utility}) on real tool-execution traces. 95\% Wilson CIs in brackets. Bold = best in each metric.}
\label{tab:main}
\begin{tabular}{llcc}
\toprule
Defense & Class & ASR $\downarrow$ [95\% CI] & Utility $\uparrow$ [95\% CI] \\
\midrule
No Defense & --- & 0.74 [0.65, 0.82] & 0.19 [0.12, 0.28] \\
\midrule
Polymorphic & IE & \textbf{0.14} [0.09, 0.23] & \textbf{0.73} [0.63, 0.81] \\
Mixture-of-Enc. & IE & 0.16 [0.10, 0.24] & 0.69 [0.59, 0.78] \\
Spotlighting & IE & 0.21 [0.14, 0.31] & 0.67 [0.57, 0.76] \\
\midrule
StruQ (prompt) & SQ & 0.70 [0.60, 0.78] & 0.17 [0.10, 0.27] \\
\midrule
FATH & RC & \textbf{0.00} [0.00, 0.04] & 0.02 [0.00, 0.07] \\
Task Shield & RC & 0.10 [0.05, 0.18] & 0.02 [0.00, 0.07] \\
IPIGuard & RC & 0.38 [0.28, 0.48] & 0.06 [0.02, 0.13] \\
\midrule
P1: Echo Canary & RC & 0.03 [0.01, 0.09] & 0.00 [0.00, 0.04] \\
P2: Orthogonal & RC & \textbf{0.00} [0.00, 0.04] & 0.00 [0.00, 0.04] \\
P3: Task Invariant & RC & 0.42 [0.32, 0.52] & 0.02 [0.00, 0.07] \\
\bottomrule
\multicolumn{4}{l}{\footnotesize IE = Input Encoding, SQ = Structured Query, RC = Runtime Checking.}
\end{tabular}
\end{table*}

\subsection{Excluded Defenses}
\label{sec:proxy_results}

Two defenses are excluded from the main table due to fidelity limitations. \textbf{Attention Tracker} (IP) requires access to model attention weights, unavailable via the GPT-4o-mini API; a token-frequency proxy would not faithfully represent the defense. \textbf{P4: Least Privilege} (AR) requires real OS-level sandboxing; a prompt-level simulation would be a category error. We discuss these qualitatively in \Cref{sec:principles} but do not report quantitative results for them.

\subsection{P2 Ablation: Does the Canary Provide Independent Security?}
\label{sec:p2_ablation}

A key question is whether P2's resistance comes from the echo-canary signal, the semantic-alignment check, or their combination. We run an ablation with three configurations (90 trials each, same tasks/attacks/seeds):

\begin{table}[t]
\centering
\small
\caption{P2 ablation: canary vs.\ alignment contributions (90 trials each).}
\label{tab:p2_ablation}
\begin{tabular}{lcc}
\toprule
Configuration & ASR $\downarrow$ & Utility $\uparrow$ \\
\midrule
Canary only & 0.03 [0.01, 0.09] & 0.00 [0.00, 0.04] \\
Alignment only & 0.09 [0.05, 0.17] & 0.04 [0.02, 0.11] \\
Combined (P2) & \textbf{0.00} [0.00, 0.04] & 0.00 [0.00, 0.04] \\
\bottomrule
\end{tabular}
\end{table}

The canary \emph{does} provide independent security value: canary-only achieves ASR 0.03 (vs.\ 0.74 baseline), a 96\% relative reduction, without the alignment check. However, all three configurations have Utility $\leq$ 0.04. The combination does not recover Utility---it merely stacks two over-blocking signals. This refutes the original hypothesis that orthogonal multi-signal combination could achieve security \emph{at acceptable utility cost}: in our end-to-end evaluation, P2's Utility is 0.00, not the 0.43 reported by the earlier proxy-based measurement. The canary's independent value is real but comes at full Utility cost.
